%============================================================
% HALAMAN SAMPUL BUKU KP (Cover Berwarna)
% Layout mengikuti Template_Buku_PA_PENS_LaTeX:
%   - Header putih: logo PENS (kiri) + "KERJA PRAKTIK" (kanan, orange)
%   - Body orange: judul kapital bold, Oleh:, mahasiswa, pembimbing, footer
% Mendukung 1, 2, atau 3 mahasiswa.
%============================================================
\begin{titlepage}
  \newgeometry{top=0pt, bottom=0pt, left=0pt, right=0pt}
  \setlength{\fboxsep}{0pt}

  % -- Bagian 1: Header Putih ----------------------------------------
  \noindent\pagecolor{white}
  \colorbox{white}{%
    \begin{minipage}[c][5.36cm][c]{\paperwidth}
      \hspace*{1.2cm}
      \begin{minipage}[t]{2.8cm}
        \vspace{0pt} % Set baseline at top
        \raisebox{-\height}{\includegraphics[width=2.5cm, keepaspectratio]{assets/pens/logo_pens.png}}
      \end{minipage}
      \hfill
      \begin{minipage}[t]{5cm}
        \vspace{0pt} % Set baseline at top
        \raggedleft
        \vspace{0.25cm}
        \raisebox{-\height}{\color{black}\FontCoverLabel\bfseries KERJA PRAKTIK}
      \end{minipage}
      \hspace*{1.2cm}
    \end{minipage}
  }\par\nointerlineskip

  % -- Bagian 2: Body Orange (latar belakang penuh) ------------------
  \noindent
  \colorbox{penskuning}{%
    \begin{minipage}[t][\dimexpr\paperheight-5.36cm\relax][t]{\paperwidth}
      \color{black}
      \begin{center}

        \vspace{0.8cm}

        % Judul: kapital, bold (TANPA italic), kata tidak dipotong
        \begin{minipage}{\LebarJudulKP}
          \centering
          \hyphenpenalty=10000
          \exhyphenpenalty=10000
          \FontCoverJudul\bfseries
          \MakeUppercase{\JudulKP}
        \end{minipage}
        \par

        \vfill

        % Blok Penulis
        {\FontCoverIsi\itshape Oleh:\par}
        \vspace{0.2cm}

        \begin{tabular}{@{}p{5.5cm}l@{}}
          \bfseries\FontCoverIsi \NamaMahasiswaSatu & \bfseries\FontCoverIsi NRP. \NRPSatu \\
          \ifcsvoid{NamaMahasiswaDua}{}{%
          \bfseries\FontCoverIsi \NamaMahasiswaDua  & \bfseries\FontCoverIsi NRP. \NRPDua  \\
          }
          \ifcsvoid{NamaMahasiswaTiga}{}{%
          \bfseries\FontCoverIsi \NamaMahasiswaTiga & \bfseries\FontCoverIsi NRP. \NRPTiga \\
          }
        \end{tabular}

        \vfill

        % Dosen Pembimbing
        {\FontCoverIsi\itshape Dosen Pembimbing:\par}
        \vspace{0.15cm}
        {\FontCoverIsi\bfseries \NamaPembimbing\par}
        {\FontCoverIsi\bfseries NIP. \NIPPembimbing\par}

        \vfill

        % Footer: institusi
        {\FontCoverIsi\bfseries \MakeUppercase{\JenjangProdi} \MakeUppercase{\NamaProdi}\par}
        {\FontCoverIsi\bfseries \MakeUppercase{\NamaDepartemen}\par}
        {\FontCoverIsi\bfseries POLITEKNIK ELEKTRONIKA NEGERI SURABAYA\par}
        {\FontCoverIsi\bfseries \TahunKP\par}
        \vspace{1.5cm}

      \end{center}
    \end{minipage}%
  }
\end{titlepage}
\restoregeometry
\pagecolor{white}
